\documentclass[12pt,a4paper]{article}

\usepackage[T2A]{fontenc}
\usepackage[utf8]{inputenc}
\usepackage[russian]{babel}
\usepackage{amsmath,amssymb,amsthm}
\usepackage{array}
\usepackage{booktabs}
\usepackage{longtable}
\usepackage{mathtools}
\usepackage{geometry}
\usepackage{enumitem}
\usepackage{tikz}
\usepackage{hyperref}
\usepackage{xcolor}
\usepackage{listings}

\geometry{margin=2.2cm}

\hypersetup{
    colorlinks=true,
    linkcolor=blue,
    urlcolor=blue
}

\lstset{
    basicstyle=\ttfamily\small,
    columns=fullflexible,
    frame=single,
    breaklines=true
}

\newtheorem{definition}{Определение}
\newtheorem{theorem}{Теорема}
\newtheorem{statement}{Утверждение}
\newtheorem{example}{Пример}

\title{Билет №23 \\ \small Процессоры общего назначения. Архитектуры CISC и RISC. Конвейер. Суперскалярность. Кэширование команд и данных. (ИТМО) \\ Конвейеризация, предсказание переходов, спекулятивное и внеочередное выполнение. Конфликты на конвейерах. Особенности конвейеризации для CISC- и RISC-процессоров. (СпБГУ)}
\author{Сериков Владислав Александрович}
\date{3 мая 2026}

\begin{document}

\maketitle
\tableofcontents
\newpage

\section{Процессоры общего назначения}

\begin{definition}
\textbf{Процессор общего назначения} --- это программно-управляемое вычислительное устройство, предназначенное для выполнения широкого класса программ, не специализированное под одну фиксированную задачу.
\end{definition}

Процессоры общего назначения используются в персональных компьютерах, серверах, мобильных устройствах, встроенных системах. Их основная особенность состоит в том, что поведение процессора задаётся не аппаратно фиксированным алгоритмом, а последовательностью машинных инструкций.

\subsection{Основные уровни описания процессора}

При изучении процессоров важно различать несколько уровней:

\begin{enumerate}
    \item \textbf{ISA} --- Instruction Set Architecture, архитектура набора команд.
    Это программно видимый интерфейс процессора:
    \begin{itemize}
        \item набор машинных инструкций;
        \item регистры общего назначения;
        \item форматы команд;
        \item способы адресации;
        \item модель памяти;
        \item исключения и прерывания;
        \item соглашения о видимом поведении команд.
    \end{itemize}

    \item \textbf{Микроархитектура} --- конкретная аппаратная реализация ISA:
    \begin{itemize}
        \item конвейер;
        \item кэши;
        \item исполнительные устройства;
        \item предсказатель переходов;
        \item буфер переупорядочивания;
        \item планировщики команд;
        \item механизмы переименования регистров.
    \end{itemize}

    \item \textbf{Физическая реализация} --- транзисторная и схемотехническая реализация:
    \begin{itemize}
        \item технология производства;
        \item тактовая частота;
        \item энергопотребление;
        \item площадь кристалла;
        \item тепловые ограничения.
    \end{itemize}
\end{enumerate}

Одна и та же ISA может иметь разные микроархитектурные реализации. Например, x86-64 может реализовываться как простыми, так и очень сложными суперскалярными процессорами с внеочередным выполнением.

\subsection{Типовая структура процессора}

Типичный процессор общего назначения содержит:

\begin{itemize}
    \item \textbf{устройство выборки команд} --- получает команды из памяти или кэша команд;
    \item \textbf{декодер команд} --- преобразует машинные команды во внутреннее представление;
    \item \textbf{регистровый файл} --- набор быстрых регистров;
    \item \textbf{арифметико-логическое устройство} --- выполняет целочисленные операции;
    \item \textbf{устройство работы с памятью} --- выполняет загрузки и сохранения;
    \item \textbf{устройство управления} --- координирует выполнение команд;
    \item \textbf{кэш-память} --- уменьшает среднее время доступа к памяти;
    \item \textbf{блок предсказания переходов} --- угадывает дальнейший поток команд;
    \item \textbf{исполнительные устройства} --- выполняют операции разных типов:
    целочисленные, вещественные, SIMD, переходы, обращения к памяти.
\end{itemize}

\section{Архитектуры CISC и RISC}

\subsection{Архитектура CISC}

\begin{definition}
\textbf{CISC} --- Complex Instruction Set Computer, архитектура с комплексным набором команд. Для CISC характерны сложные команды, большое количество форматов инструкций и развитые способы адресации.
\end{definition}

Классический пример CISC-архитектуры --- x86.

\subsubsection{Основные признаки CISC}

Для CISC обычно характерны:

\begin{itemize}
    \item команды переменной длины;
    \item большое количество инструкций;
    \item наличие сложных инструкций, выполняющих несколько элементарных действий;
    \item операции, которые могут напрямую работать с памятью;
    \item большое число способов адресации;
    \item сложный декодер;
    \item историческая ориентация на компактность машинного кода.
\end{itemize}

Например, в CISC возможна команда вида:

\[
    \texttt{ADD [addr], R1}
\]

которая может означать:

\begin{enumerate}
    \item прочитать значение из памяти по адресу \texttt{addr};
    \item прибавить к нему значение регистра \texttt{R1};
    \item записать результат обратно в память.
\end{enumerate}

То есть одна машинная команда может включать несколько микроопераций.

\subsubsection{Преимущества CISC}

\begin{itemize}
    \item высокая плотность кода;
    \item удобство для компиляторов ранних поколений и ассемблерного программирования;
    \item богатые способы адресации;
    \item обратная совместимость, особенно важная для массовых платформ.
\end{itemize}

\subsubsection{Недостатки CISC}

\begin{itemize}
    \item сложный декодер команд;
    \item неоднородное время выполнения инструкций;
    \item трудности глубокой конвейеризации;
    \item переменная длина команд усложняет выборку и декодирование;
    \item сложность аппаратной реализации.
\end{itemize}

\subsection{Архитектура RISC}

\begin{definition}
\textbf{RISC} --- Reduced Instruction Set Computer, архитектура с сокращённым набором команд. Для RISC характерны простые команды, регулярные форматы инструкций и ориентация на эффективную конвейеризацию.
\end{definition}

Примеры RISC-архитектур:

\begin{itemize}
    \item MIPS;
    \item ARM;
    \item SPARC;
    \item PowerPC;
    \item RISC-V.
\end{itemize}

\subsubsection{Основные признаки RISC}

Для RISC обычно характерны:

\begin{itemize}
    \item фиксированная или почти фиксированная длина команд;
    \item малое количество простых форматов инструкций;
    \item архитектура типа \textbf{load-store};
    \item арифметические операции выполняются только над регистрами;
    \item доступ к памяти осуществляется специальными командами загрузки и сохранения;
    \item большое число регистров общего назначения;
    \item простое и быстрое декодирование;
    \item удобство конвейеризации.
\end{itemize}

\begin{definition}
\textbf{Load-store архитектура} --- это архитектура, в которой обращения к памяти выполняются только специальными командами загрузки и сохранения, а арифметико-логические операции работают только с регистрами.
\end{definition}

Например, операция

\[
    A = A + B
\]

в RISC-стиле может быть представлена так:

\begin{lstlisting}
LOAD R1, [A]     ; R1 = Memory[A]
LOAD R2, [B]     ; R2 = Memory[B]
ADD  R3, R1, R2  ; R3 = R1 + R2
STORE [A], R3    ; Memory[A] = R3
\end{lstlisting}

\subsubsection{Преимущества RISC}

\begin{itemize}
    \item простая аппаратная реализация;
    \item удобство конвейеризации;
    \item регулярные форматы команд;
    \item высокая тактовая частота;
    \item эффективная работа компилятора;
    \item проще реализовать суперскалярность и внеочередное выполнение.
\end{itemize}

\subsubsection{Недостатки RISC}

\begin{itemize}
    \item меньшая плотность кода по сравнению с CISC;
    \item большее число инструкций для некоторых операций;
    \item большая зависимость производительности от качества компилятора;
    \item необходимость большего регистрового файла.
\end{itemize}

\subsection{Сравнение CISC и RISC}

\begin{longtable}{p{0.27\textwidth}p{0.34\textwidth}p{0.34\textwidth}}
\toprule
\textbf{Признак} & \textbf{CISC} & \textbf{RISC} \\
\midrule
Длина команды & Обычно переменная & Обычно фиксированная \\
\midrule
Сложность команд & Высокая & Низкая \\
\midrule
Форматы команд & Много форматов & Мало форматов \\
\midrule
Работа с памятью & Многие команды могут обращаться к памяти & Только load/store \\
\midrule
Декодирование & Сложное & Простое \\
\midrule
Плотность кода & Обычно выше & Обычно ниже \\
\midrule
Конвейеризация & Сложнее & Проще \\
\midrule
Пример & x86, x86-64 & ARM, MIPS, RISC-V \\
\bottomrule
\end{longtable}

\subsection{Современное сближение CISC и RISC}

Современные CISC-процессоры часто используют внутреннее RISC-подобное представление. Например, сложные x86-команды декодируются во внутренние микрооперации.

\begin{definition}
\textbf{Микрооперация} --- это простая внутренняя операция процессора, на которую может быть разложена одна машинная инструкция.
\end{definition}

Схематически:

\[
    \text{x86-инструкция}
    \xrightarrow{\text{декодер}}
    \text{микрооперации}
    \xrightarrow{\text{планировщик}}
    \text{исполнительные устройства}.
\]

Таким образом, внешняя ISA может быть CISC, а внутренняя микроархитектура --- RISC-подобной.

\section{Производительность процессора}

\subsection{Основная формула времени выполнения}

Время выполнения программы можно представить как:

\[
    T = IC \cdot CPI \cdot T_{\text{cycle}},
\]

где:

\begin{itemize}
    \item $T$ --- время выполнения программы;
    \item $IC$ --- instruction count, количество выполненных машинных инструкций;
    \item $CPI$ --- cycles per instruction, среднее число тактов на инструкцию;
    \item $T_{\text{cycle}}$ --- длительность одного такта.
\end{itemize}

Так как

\[
    f = \frac{1}{T_{\text{cycle}}},
\]

то:

\[
    T = \frac{IC \cdot CPI}{f}.
\]

\subsection{Интерпретация формулы}

Для ускорения программы можно:

\begin{enumerate}
    \item уменьшить число инструкций $IC$;
    \item уменьшить среднее число тактов на инструкцию $CPI$;
    \item увеличить тактовую частоту $f$.
\end{enumerate}

Но эти параметры связаны между собой. Например, более сложная ISA может уменьшать $IC$, но увеличивать $CPI$ или снижать максимальную частоту.

\section{Конвейеризация}

\subsection{Идея конвейеризации}

\begin{definition}
\textbf{Конвейеризация} --- это способ организации выполнения команд, при котором выполнение каждой команды разбивается на несколько стадий, а разные команды одновременно находятся на разных стадиях.
\end{definition}

Аналогия: на заводском конвейере одно изделие красят, другое собирают, третье упаковывают. Каждое изделие проходит все стадии, но одновременно обрабатываются разные изделия.

\subsection{Классический пятистадийный RISC-конвейер}

Классическая учебная модель RISC-конвейера содержит пять стадий:

\begin{enumerate}
    \item \textbf{IF} --- Instruction Fetch, выборка инструкции;
    \item \textbf{ID} --- Instruction Decode, декодирование и чтение регистров;
    \item \textbf{EX} --- Execute, выполнение операции или вычисление адреса;
    \item \textbf{MEM} --- Memory Access, обращение к памяти;
    \item \textbf{WB} --- Write Back, запись результата в регистр.
\end{enumerate}

\[
\begin{array}{c|ccccc}
\text{Такт} & 1 & 2 & 3 & 4 & 5 \\
\hline
I_1 & IF & ID & EX & MEM & WB
\end{array}
\]

При последовательном, неконвейерном выполнении следующая инструкция началась бы только после завершения предыдущей. В конвейере следующая инструкция начинает выполняться уже на следующем такте.

\[
\begin{array}{c|ccccccccc}
\text{Такт} & 1 & 2 & 3 & 4 & 5 & 6 & 7 & 8 & 9 \\
\hline
I_1 & IF & ID & EX & MEM & WB &   &   &   &   \\
I_2 &    & IF & ID & EX & MEM & WB &   &   &   \\
I_3 &    &    & IF & ID & EX & MEM & WB &   &   \\
I_4 &    &    &    & IF & ID & EX & MEM & WB &   \\
I_5 &    &    &    &    & IF & ID & EX & MEM & WB
\end{array}
\]

\subsection{Ускорение от конвейеризации}

Пусть выполнение одной инструкции разбито на $k$ стадий, каждая стадия занимает один такт. Тогда:

\begin{itemize}
    \item без конвейера $n$ инструкций выполняются примерно за $nk$ тактов;
    \item с конвейером $n$ инструкций выполняются примерно за $k+n-1$ тактов.
\end{itemize}

\begin{theorem}
В идеальном $k$-стадийном конвейере без конфликтов и простоев ускорение при выполнении $n$ инструкций равно
\[
    S(n) = \frac{nk}{k+n-1}.
\]
При $n \to \infty$ ускорение стремится к $k$.
\end{theorem}

\begin{proof}
Без конвейеризации каждая инструкция проходит $k$ стадий последовательно. Поэтому для $n$ инструкций требуется
\[
    T_{\text{seq}} = nk
\]
тактов.

В конвейере первая инструкция завершится через $k$ тактов. После заполнения конвейера на каждом следующем такте завершается одна инструкция. Следовательно, для завершения всех $n$ инструкций требуется:
\[
    T_{\text{pipe}} = k + (n-1) = k+n-1.
\]

Ускорение равно отношению времени без конвейера ко времени с конвейером:
\[
    S(n) = \frac{T_{\text{seq}}}{T_{\text{pipe}}}
         = \frac{nk}{k+n-1}.
\]

Теперь найдём предел:
\[
    \lim_{n\to\infty} \frac{nk}{k+n-1}
    =
    \lim_{n\to\infty} \frac{k}{1+\frac{k-1}{n}}
    =
    k.
\]

Значит, максимальное идеальное ускорение стремится к числу стадий конвейера.
\end{proof}

\subsection{Минимальный пример}

Пусть $k=5$, $n=10$.

Без конвейера:

\[
    T_{\text{seq}} = nk = 10 \cdot 5 = 50.
\]

С конвейером:

\[
    T_{\text{pipe}} = k+n-1 = 5+10-1=14.
\]

Ускорение:

\[
    S = \frac{50}{14} \approx 3.57.
\]

Хотя конвейер имеет 5 стадий, ускорение меньше 5, потому что нужно время на заполнение и опустошение конвейера.

\subsection{Идеальный CPI конвейера}

В идеальном скалярном конвейере после заполнения завершается одна инструкция за такт. Поэтому:

\[
    CPI_{\text{ideal}} = 1.
\]

На практике:

\[
    CPI = 1 + \text{штрафы за конфликты}.
\]

\section{Конфликты в конвейерах}

\begin{definition}
\textbf{Конфликт конвейера} --- это ситуация, при которой следующая инструкция не может быть выполнена в запланированном такте без нарушения корректности или из-за нехватки аппаратного ресурса.
\end{definition}

Основные классы конфликтов:

\begin{enumerate}
    \item структурные;
    \item по данным;
    \item управляющие.
\end{enumerate}

\subsection{Структурные конфликты}

\begin{definition}
\textbf{Структурный конфликт} возникает, когда нескольким стадиям конвейера одновременно требуется один и тот же аппаратный ресурс.
\end{definition}

Пример: если команды и данные хранятся в единой памяти с одним портом, то в один такт процессор может захотеть:

\begin{itemize}
    \item выбрать следующую инструкцию;
    \item выполнить загрузку данных.
\end{itemize}

Обе операции требуют доступа к памяти, возникает конфликт.

\subsubsection{Способы устранения}

\begin{itemize}
    \item разделение кэша команд и кэша данных;
    \item многопортовые регистровые файлы;
    \item дублирование исполнительных устройств;
    \item введение очередей и буферов;
    \item остановка конвейера при невозможности устранить конфликт.
\end{itemize}

\subsection{Конфликты по данным}

\begin{definition}
\textbf{Конфликт по данным} возникает, когда одна инструкция зависит от результата другой инструкции, ещё не завершившей выполнение.
\end{definition}

Различают три типа зависимостей.

\subsubsection{RAW: Read After Write}

\begin{definition}
\textbf{RAW-зависимость} --- истинная зависимость по данным: инструкция читает значение, которое должна записать предыдущая инструкция.
\end{definition}

Пример:

\begin{lstlisting}
ADD R1, R2, R3   ; R1 = R2 + R3
SUB R4, R1, R5   ; R4 = R1 - R5
\end{lstlisting}

Инструкция \texttt{SUB} должна использовать новое значение \texttt{R1}, вычисленное командой \texttt{ADD}.

\subsubsection{WAR: Write After Read}

\begin{definition}
\textbf{WAR-зависимость} --- антизависимость: более поздняя инструкция записывает регистр, который более ранняя инструкция должна успеть прочитать.
\end{definition}

Пример:

\begin{lstlisting}
ADD R4, R1, R2   ; reads R1
MOV R1, R5       ; writes R1
\end{lstlisting}

Если \texttt{MOV} выполнится слишком рано, \texttt{ADD} может прочитать неправильное значение.

В простом пятистадийном конвейере с последовательной выдачей WAR обычно не возникает, потому что чтение регистров происходит раньше записи более поздней инструкции. Но при внеочередном выполнении WAR становится важной проблемой.

\subsubsection{WAW: Write After Write}

\begin{definition}
\textbf{WAW-зависимость} --- выходная зависимость: две инструкции записывают один и тот же регистр, и важно сохранить правильный порядок записей.
\end{definition}

Пример:

\begin{lstlisting}
MUL R1, R2, R3   ; R1 = R2 * R3
ADD R1, R4, R5   ; R1 = R4 + R5
\end{lstlisting}

Итоговое значение \texttt{R1} должно соответствовать более поздней инструкции \texttt{ADD}. Если \texttt{MUL} завершится позже и перезапишет результат, возникнет ошибка.

\subsection{Методы борьбы с конфликтами по данным}

\subsubsection{Остановка конвейера}

Самый простой метод --- вставить пустые такты.

\begin{definition}
\textbf{Пузырь конвейера} --- это пустая операция, вставленная в конвейер для ожидания готовности данных или ресурса.
\end{definition}

Пример:

\begin{lstlisting}
ADD R1, R2, R3
NOP
NOP
SUB R4, R1, R5
\end{lstlisting}

Команды \texttt{NOP} ничего не делают, но дают время первой команде записать результат.

\subsubsection{Forwarding}

\begin{definition}
\textbf{Forwarding}, или \textbf{bypassing}, --- это передача результата напрямую из промежуточной стадии конвейера на вход исполнительного устройства, минуя запись и последующее чтение из регистрового файла.
\end{definition}

Пример:

\begin{lstlisting}
ADD R1, R2, R3
SUB R4, R1, R5
\end{lstlisting}

Результат \texttt{ADD} может быть готов уже в конце стадии \texttt{EX}. Его можно передать напрямую на вход \texttt{SUB}, которая тоже находится на стадии \texttt{EX} в следующем такте.

Схема:

\[
\begin{array}{c|ccccc}
\text{Такт} & 1 & 2 & 3 & 4 & 5 \\
\hline
ADD & IF & ID & EX & MEM & WB \\
SUB &    & IF & ID & EX  & MEM
\end{array}
\]

Результат \texttt{ADD} передаётся из \texttt{EX/MEM} в \texttt{EX} команды \texttt{SUB}.

\subsubsection{Перестановка инструкций компилятором}

Компилятор может расположить независимые инструкции между зависимыми.

Было:

\begin{lstlisting}
LOAD R1, [A]
ADD  R2, R1, R3
MUL  R4, R5, R6
\end{lstlisting}

Стало:

\begin{lstlisting}
LOAD R1, [A]
MUL  R4, R5, R6
ADD  R2, R1, R3
\end{lstlisting}

Если \texttt{MUL} не зависит от \texttt{LOAD}, она заполняет потенциальный простой.

\subsubsection{Переименование регистров}

\begin{definition}
\textbf{Переименование регистров} --- это микроархитектурный метод, при котором архитектурные регистры отображаются на большее множество физических регистров для устранения ложных зависимостей WAR и WAW.
\end{definition}

Пример WAW:

\begin{lstlisting}
MUL R1, R2, R3
ADD R1, R4, R5
\end{lstlisting}

После переименования:

\begin{lstlisting}
MUL P10, P2, P3
ADD P11, P4, P5
\end{lstlisting}

Обе инструкции пишут в разные физические регистры, поэтому конфликт WAW исчезает. Архитектурный регистр \texttt{R1} после завершения второй инструкции будет связан с \texttt{P11}.

\subsection{Управляющие конфликты}

\begin{definition}
\textbf{Управляющий конфликт} возникает из-за команд изменения потока управления: условных и безусловных переходов, вызовов функций, возвратов, исключений.
\end{definition}

Пример:

\begin{lstlisting}
BEQ R1, R2, label
ADD R3, R4, R5
SUB R6, R7, R8
label:
MUL R9, R10, R11
\end{lstlisting}

Пока не известно, выполнится ли переход, не ясно, какую инструкцию выбирать дальше: \texttt{ADD} или \texttt{MUL}.

\subsubsection{Методы борьбы}

\begin{itemize}
    \item остановка до разрешения перехода;
    \item вычисление условия перехода как можно раньше;
    \item статическое предсказание;
    \item динамическое предсказание;
    \item спекулятивное выполнение;
    \item delayed branch в некоторых RISC-архитектурах.
\end{itemize}

\section{Предсказание переходов}

\subsection{Зачем нужно предсказание}

В глубоком конвейере процессор должен выбирать инструкции заранее. Если встречается условный переход, то до вычисления условия неизвестно, какой адрес будет следующим.

Без предсказания процессор вынужден ждать, что приводит к потере тактов.

\begin{definition}
\textbf{Предсказание переходов} --- это аппаратный или программный механизм, который пытается заранее определить, будет ли выполнен условный переход и каков будет следующий адрес выборки команд.
\end{definition}

\subsection{Два аспекта предсказания}

Для перехода нужно предсказать:

\begin{enumerate}
    \item \textbf{направление}: переход будет выполнен или нет;
    \item \textbf{целевой адрес}: откуда продолжать выборку команд.
\end{enumerate}

Для второго обычно используется \textbf{BTB}.

\begin{definition}
\textbf{BTB} --- Branch Target Buffer, буфер целевых адресов переходов. Это кэш, в котором для ранее встречавшихся переходов хранятся адреса их целей.
\end{definition}

\subsection{Статическое предсказание}

\begin{definition}
\textbf{Статическое предсказание} --- это предсказание, не использующее историю выполнения программы.
\end{definition}

Примеры:

\begin{itemize}
    \item всегда считать, что переход не выполняется;
    \item всегда считать, что переход выполняется;
    \item обратные переходы считать выполняющимися, прямые --- нет;
    \item использовать подсказки компилятора.
\end{itemize}

Эвристика ``обратный переход выполняется'' хорошо работает для циклов.

Пример:

\begin{lstlisting}
loop:
    ...
    SUB R1, R1, #1
    BNE R1, R0, loop
\end{lstlisting}

Переход назад к \texttt{loop} обычно выполняется много раз, пока цикл не завершится.

\subsection{Динамическое предсказание}

\begin{definition}
\textbf{Динамическое предсказание} --- это предсказание, использующее информацию о предыдущем поведении переходов во время выполнения программы.
\end{definition}

\subsubsection{Однобитный предсказатель}

Для каждого перехода хранится один бит:

\[
    b =
    \begin{cases}
    1, & \text{предсказывать taken};\\
    0, & \text{предсказывать not taken}.
    \end{cases}
\]

После фактического выполнения перехода бит обновляется на реальный исход.

\paragraph{Алгоритм однобитного предсказателя}

\begin{enumerate}
    \item При выборке команды перехода найти её запись в таблице истории.
    \item Если бит равен $1$, предсказать ``переход выполнится''.
    \item Если бит равен $0$, предсказать ``переход не выполнится''.
    \item После вычисления реального исхода записать в бит новое значение.
\end{enumerate}

\paragraph{Минимальный пример}

Пусть цикл выполняется 5 раз, затем завершается. Последовательность исходов:

\[
    T, T, T, T, N.
\]

Если изначально предсказывается $T$, то первые четыре перехода предсказаны верно, последний --- неверно.

Проблема однобитного предсказателя: для циклов он часто ошибается на выходе из цикла и затем может ошибиться на первом входе в цикл при следующем запуске.

\subsubsection{Двухбитный насыщаемый счётчик}

Более устойчивый вариант --- двухбитный предсказатель.

Состояния:

\[
\begin{array}{c|c|c}
\text{Состояние} & \text{Название} & \text{Предсказание} \\
\hline
00 & \text{сильно not taken} & N \\
01 & \text{слабо not taken} & N \\
10 & \text{слабо taken} & T \\
11 & \text{сильно taken} & T
\end{array}
\]

При фактическом $T$ счётчик увеличивается на 1, но не выше $3$.
При фактическом $N$ счётчик уменьшается на 1, но не ниже $0$.

\paragraph{Алгоритм двухбитного предсказателя}

Пусть $c \in \{0,1,2,3\}$.

\begin{enumerate}
    \item Если $c \geq 2$, предсказать $T$.
    \item Если $c < 2$, предсказать $N$.
    \item После реального исхода:
    \begin{itemize}
        \item если исход $T$, заменить $c$ на $\min(c+1,3)$;
        \item если исход $N$, заменить $c$ на $\max(c-1,0)$.
    \end{itemize}
\end{enumerate}

\paragraph{Иллюстрация автомата}

\[
\begin{array}{ccccc}
00 & \xleftrightarrow{\ T/N\ } & 01 & \xleftrightarrow{\ T/N\ } & 10 \xleftrightarrow{\ T/N\ } 11
\end{array}
\]

Более точно:

\[
00 \xrightarrow{T} 01 \xrightarrow{T} 10 \xrightarrow{T} 11,
\]
\[
11 \xrightarrow{N} 10 \xrightarrow{N} 01 \xrightarrow{N} 00.
\]

\paragraph{Минимальный пример}

Пусть начальное состояние $11$, а исходы цикла:

\[
    T,T,T,T,N.
\]

Таблица:

\[
\begin{array}{c|c|c|c}
\text{Шаг} & \text{Состояние до} & \text{Предсказание} & \text{Реальный исход} \\
\hline
1 & 11 & T & T \\
2 & 11 & T & T \\
3 & 11 & T & T \\
4 & 11 & T & T \\
5 & 11 & T & N
\end{array}
\]

После последнего шага состояние станет $10$, но предсказание останется $T$.
То есть один нетипичный исход не меняет направление предсказания сразу.

\subsection{Глобальная и локальная история}

\begin{definition}
\textbf{Локальная история перехода} --- история исходов одного конкретного перехода.
\end{definition}

\begin{definition}
\textbf{Глобальная история} --- история исходов последних переходов независимо от того, какими командами они были вызваны.
\end{definition}

Глобальная история позволяет предсказывать коррелированные переходы.

Пример:

\begin{lstlisting}
if (x < 0)  ...
if (x >= 0) ...
\end{lstlisting}

Если первый переход выполнен, исход второго может быть предсказан с учётом первого.

\section{Спекулятивное выполнение}

\begin{definition}
\textbf{Спекулятивное выполнение} --- это выполнение инструкций до того, как точно известно, понадобятся ли их результаты в архитектурном состоянии программы.
\end{definition}

Обычно спекуляция связана с предсказанием переходов. Если процессор предсказал, что переход будет выполнен, он начинает выполнять команды по предсказанному адресу.

\subsection{Корректная спекуляция}

Если предсказание верно, результаты спекулятивных инструкций можно зафиксировать.

Если предсказание неверно:

\begin{enumerate}
    \item спекулятивные инструкции отменяются;
    \item конвейер очищается;
    \item выборка продолжается с правильного адреса;
    \item процессор теряет несколько тактов.
\end{enumerate}

\begin{definition}
\textbf{Штраф за ошибку предсказания} --- число тактов, потерянных из-за неверного предсказания перехода.
\end{definition}

\subsection{Точное исключение}

Спекуляция требует аккуратного обращения с исключениями.

\begin{definition}
\textbf{Точное исключение} --- это исключение, при котором архитектурное состояние процессора выглядит так, будто все инструкции до исключившейся завершены, а сама исключившаяся и все последующие инструкции не изменили архитектурное состояние.
\end{definition}

Для поддержки точных исключений современные процессоры используют буфер переупорядочивания.

\section{Суперскалярность}

\begin{definition}
\textbf{Суперскалярный процессор} --- это процессор, способный выдавать, выполнять или завершать более одной инструкции за такт.
\end{definition}

Если обычный скалярный конвейер в идеале завершает одну инструкцию за такт, то суперскалярный процессор ширины $w$ в идеале может завершать до $w$ инструкций за такт.

\subsection{Ширина суперскалярного процессора}

\begin{definition}
\textbf{Ширина процессора} --- максимальное число инструкций или микроопераций, которые процессор может обработать на некоторой стадии за один такт.
\end{definition}

Различают:

\begin{itemize}
    \item ширину выборки;
    \item ширину декодирования;
    \item ширину выдачи;
    \item ширину исполнения;
    \item ширину завершения.
\end{itemize}

Например, если процессор может декодировать 4 инструкции за такт, но завершать только 2, то узким местом может стать стадия завершения.

\subsection{ILP}

\begin{definition}
\textbf{ILP} --- Instruction-Level Parallelism, параллелизм уровня инструкций. Это возможность выполнять несколько инструкций одной программы параллельно или с перекрытием.
\end{definition}

Пример независимых инструкций:

\begin{lstlisting}
ADD R1, R2, R3
SUB R4, R5, R6
MUL R7, R8, R9
\end{lstlisting}

Эти команды не зависят друг от друга и потенциально могут выполняться одновременно.

Пример зависимых инструкций:

\begin{lstlisting}
ADD R1, R2, R3
SUB R4, R1, R5
MUL R6, R4, R7
\end{lstlisting}

Здесь каждая следующая команда зависит от результата предыдущей, поэтому параллелизм ограничен.

\subsection{Теоретический предел IPC}

\begin{definition}
\textbf{IPC} --- Instructions Per Cycle, среднее число инструкций, завершаемых за один такт.
\end{definition}

\[
    IPC = \frac{1}{CPI}.
\]

Для суперскалярного процессора ширины $w$:

\[
    IPC \leq w.
\]

На практике IPC меньше $w$ из-за:

\begin{itemize}
    \item зависимостей по данным;
    \item промахов кэша;
    \item ошибок предсказания переходов;
    \item ограничений исполнительных устройств;
    \item ограниченной ширины декодирования;
    \item недостаточного параллелизма в программе.
\end{itemize}

\section{Внеочередное выполнение}

\subsection{Идея внеочередного выполнения}

\begin{definition}
\textbf{Внеочередное выполнение} --- это способ выполнения, при котором процессор может выполнять инструкции не в программном порядке, если это не нарушает видимого результата программы.
\end{definition}

Главная цель --- использовать простаивающие исполнительные устройства, пока некоторые инструкции ожидают данные.

Пример:

\begin{lstlisting}
LOAD R1, [A]       ; long memory wait
ADD  R2, R1, R3    ; depends on LOAD
MUL  R4, R5, R6    ; independent
SUB  R7, R8, R9    ; independent
\end{lstlisting}

Если \texttt{LOAD} задерживается из-за промаха кэша, команды \texttt{MUL} и \texttt{SUB} можно выполнить раньше \texttt{ADD}, потому что они не зависят от \texttt{LOAD}.

\subsection{Основные механизмы}

Для внеочередного выполнения нужны:

\begin{itemize}
    \item переименование регистров;
    \item очередь инструкций или станции резервирования;
    \item отслеживание готовности операндов;
    \item динамическое планирование;
    \item буфер переупорядочивания;
    \item механизм завершения в программном порядке.
\end{itemize}

\subsection{Алгоритм Томасуло}

Классический алгоритм динамического планирования --- алгоритм Томасуло.

\begin{definition}
\textbf{Алгоритм Томасуло} --- это аппаратный алгоритм динамического планирования инструкций, использующий станции резервирования, переименование регистров и рассылку результатов по общей шине данных.
\end{definition}

\subsubsection{Основные структуры}

\begin{itemize}
    \item \textbf{станции резервирования} --- буферы около исполнительных устройств;
    \item \textbf{таблица состояния регистров} --- указывает, какая станция должна произвести значение регистра;
    \item \textbf{общая шина данных} --- используется для рассылки результатов ожидающим инструкциям;
    \item \textbf{исполнительные устройства} --- выполняют операции.
\end{itemize}

\subsubsection{Шаги алгоритма}

Алгоритм обычно описывается тремя фазами.

\paragraph{1. Issue --- выдача инструкции}

\begin{enumerate}
    \item Декодировать очередную инструкцию.
    \item Найти свободную станцию резервирования нужного типа.
    \item Если станции нет, остановить выдачу.
    \item Если операнд уже готов, записать его значение в станцию.
    \item Если операнд не готов, записать тег станции, которая его произведёт.
    \item Для результата обновить таблицу состояния регистров: теперь соответствующий регистр ожидает результат от данной станции.
\end{enumerate}

\paragraph{2. Execute --- выполнение}

\begin{enumerate}
    \item Станция резервирования ждёт готовности всех операндов.
    \item Когда операнды готовы и исполнительное устройство свободно, инструкция начинает выполнение.
    \item Операция выполняется требуемое число тактов.
\end{enumerate}

\paragraph{3. Write result --- запись результата}

\begin{enumerate}
    \item Результат помещается на общую шину данных.
    \item Все станции резервирования, ожидающие этот тег, получают значение.
    \item Если таблица состояния регистра всё ещё указывает на данную станцию, регистр получает значение.
    \item Станция освобождается.
\end{enumerate}

\subsubsection{Минимальный пример}

Пусть есть инструкции:

\begin{lstlisting}
I1: LOAD F2, 0(R1)
I2: MUL  F4, F2, F6
I3: ADD  F8, F10, F12
\end{lstlisting}

Зависимости:

\begin{itemize}
    \item \texttt{I2} зависит от \texttt{I1}, потому что использует \texttt{F2};
    \item \texttt{I3} не зависит от \texttt{I1} и \texttt{I2}.
\end{itemize}

При внеочередном выполнении:

\begin{enumerate}
    \item \texttt{I1} выдаётся и начинает загрузку.
    \item \texttt{I2} выдаётся, но ждёт результат \texttt{F2}.
    \item \texttt{I3} выдаётся, её операнды готовы.
    \item \texttt{I3} может выполниться раньше \texttt{I2}.
    \item После завершения \texttt{LOAD} результат рассылается, и \texttt{I2} начинает выполнение.
\end{enumerate}

\subsection{Буфер переупорядочивания}

\begin{definition}
\textbf{ROB} --- Reorder Buffer, буфер переупорядочивания. Это структура, которая хранит информацию о выполняющихся инструкциях и позволяет фиксировать их результаты в программном порядке.
\end{definition}

ROB нужен для:

\begin{itemize}
    \item точных исключений;
    \item отмены спекулятивных инструкций;
    \item завершения команд в программном порядке;
    \item восстановления состояния после ошибки предсказания.
\end{itemize}

\subsection{Commit}

\begin{definition}
\textbf{Commit}, или retirement, --- стадия окончательной фиксации результата инструкции в архитектурном состоянии процессора.
\end{definition}

Инструкция может быть выполнена вне очереди, но зафиксирована обычно должна быть в программном порядке.

\section{Кэширование команд и данных}

\subsection{Проблема разрыва между процессором и памятью}

Современный процессор работает значительно быстрее основной памяти. Если бы каждая команда и каждое данное брались напрямую из DRAM, процессор большую часть времени простаивал бы.

Для решения используется иерархия памяти:

\[
\text{регистры}
\rightarrow
L1
\rightarrow
L2
\rightarrow
L3
\rightarrow
DRAM
\rightarrow
SSD/HDD.
\]

Чем ближе память к процессору, тем она быстрее, меньше и дороже в расчёте на байт.

\subsection{Принцип локальности}

Кэширование основано на принципе локальности.

\begin{definition}
\textbf{Временная локальность} --- если некоторый объект памяти был использован, вероятно, он скоро будет использован снова.
\end{definition}

\begin{definition}
\textbf{Пространственная локальность} --- если был использован некоторый адрес, вероятно, скоро будут использованы соседние адреса.
\end{definition}

Пример временной локальности:

\begin{lstlisting}
for (int i = 0; i < n; ++i)
    sum += a[i];
\end{lstlisting}

Переменная \texttt{sum} используется много раз.

Пример пространственной локальности: последовательный обход массива \texttt{a[i]}.

\subsection{Кэш-линия}

\begin{definition}
\textbf{Кэш-линия} --- минимальный блок данных, который переносится между уровнями памяти и кэша.
\end{definition}

Например, если размер кэш-линии равен 64 байтам, то при обращении к одному байту из памяти в кэш обычно загружается весь блок из 64 байт.

\subsection{Попадания и промахи}

\begin{definition}
\textbf{Кэш-попадание} --- ситуация, когда требуемый блок найден в кэше.
\end{definition}

\begin{definition}
\textbf{Кэш-промах} --- ситуация, когда требуемого блока нет в кэше и его нужно загрузить с более медленного уровня памяти.
\end{definition}

Среднее время доступа к памяти:

\[
    AMAT = T_{\text{hit}} + R_{\text{miss}} \cdot P_{\text{miss}},
\]

где:

\begin{itemize}
    \item $AMAT$ --- average memory access time;
    \item $T_{\text{hit}}$ --- время доступа при попадании;
    \item $R_{\text{miss}}$ --- доля промахов;
    \item $P_{\text{miss}}$ --- штраф за промах.
\end{itemize}

\subsection{Кэш команд и кэш данных}

В процессорах часто используется раздельный L1-кэш:

\begin{itemize}
    \item \textbf{I-cache} --- instruction cache, кэш команд;
    \item \textbf{D-cache} --- data cache, кэш данных.
\end{itemize}

Разделение полезно потому, что в одном такте процессору может понадобиться:

\begin{enumerate}
    \item выбрать следующую инструкцию;
    \item выполнить загрузку или сохранение данных.
\end{enumerate}

Если кэш единый и однопортовый, возникнет структурный конфликт. Раздельные I-cache и D-cache позволяют выполнять эти операции параллельно.

\subsection{Организация кэша}

Кэш обычно делится на наборы.

\begin{definition}
\textbf{Прямо отображаемый кэш} --- кэш, в котором каждый блок памяти может попасть только в одно строго определённое место кэша.
\end{definition}

\begin{definition}
\textbf{Полностью ассоциативный кэш} --- кэш, в котором блок памяти может быть размещён в любой строке кэша.
\end{definition}

\begin{definition}
\textbf{Множественно-ассоциативный кэш} --- кэш, в котором блок памяти может быть размещён в одном из нескольких мест внутри соответствующего набора.
\end{definition}

Адрес обычно разбивается на поля:

\[
    \text{адрес} =
    \text{tag} \mid \text{index} \mid \text{offset}.
\]

\begin{itemize}
    \item \textbf{offset} выбирает байт внутри кэш-линии;
    \item \textbf{index} выбирает набор;
    \item \textbf{tag} проверяет, тот ли блок находится в выбранной строке.
\end{itemize}

\subsection{Пример разбиения адреса}

Пусть:

\begin{itemize}
    \item адрес 32-битный;
    \item размер кэша $32$ КБ;
    \item размер строки $64$ байта;
    \item кэш 4-ассоциативный.
\end{itemize}

Число строк:

\[
    \frac{32 \cdot 1024}{64} = 512.
\]

Число наборов:

\[
    \frac{512}{4} = 128.
\]

Для смещения внутри строки нужно:

\[
    \log_2 64 = 6
\]

бит.

Для индекса набора нужно:

\[
    \log_2 128 = 7
\]

бит.

Для тега остаётся:

\[
    32 - 6 - 7 = 19
\]

бит.

Итого:

\[
    \underbrace{19}_{tag}
    \mid
    \underbrace{7}_{index}
    \mid
    \underbrace{6}_{offset}.
\]

\subsection{Политики записи}

Для D-cache важны политики записи.

\begin{definition}
\textbf{Write-through} --- политика, при которой запись одновременно выполняется в кэш и на следующий уровень памяти.
\end{definition}

\begin{definition}
\textbf{Write-back} --- политика, при которой запись сначала выполняется только в кэш, а на следующий уровень памяти изменённая строка записывается при вытеснении.
\end{definition}

Write-through проще, но создаёт больше трафика.
Write-back эффективнее, но требует хранить dirty-бит.

\begin{definition}
\textbf{Dirty-бит} --- признак того, что кэш-линия изменена относительно копии на нижнем уровне памяти.
\end{definition}

\subsection{Типы кэш-промахов}

Классическая классификация: три C.

\begin{enumerate}
    \item \textbf{Compulsory misses} --- обязательные промахи первого обращения.
    \item \textbf{Capacity misses} --- промахи из-за недостаточного размера кэша.
    \item \textbf{Conflict misses} --- промахи из-за конкуренции блоков за один набор.
\end{enumerate}

\section{Особенности конвейеризации RISC- и CISC-процессоров}

\subsection{RISC и конвейеризация}

RISC-архитектуры исторически хорошо подходят для конвейеризации:

\begin{itemize}
    \item команды имеют регулярный формат;
    \item декодирование простое;
    \item доступ к памяти отделён от арифметических операций;
    \item стадии конвейера легче сделать примерно равными по длительности;
    \item проще определить зависимости;
    \item команды часто выполняют одну простую операцию.
\end{itemize}

Классический пятистадийный конвейер хорошо соответствует RISC load-store модели:

\[
    IF \rightarrow ID \rightarrow EX \rightarrow MEM \rightarrow WB.
\]

\subsection{CISC и конвейеризация}

CISC усложняет конвейеризацию по нескольким причинам:

\begin{itemize}
    \item команды переменной длины;
    \item границы инструкций не всегда легко определить заранее;
    \item одна команда может выполнять много действий;
    \item команда может обращаться к памяти несколько раз;
    \item время выполнения команд сильно различается;
    \item сложные способы адресации требуют дополнительной логики.
\end{itemize}

Например, CISC-команда может:

\begin{enumerate}
    \item прочитать операнд из памяти;
    \item выполнить сложное вычисление адреса;
    \item выполнить арифметическую операцию;
    \item записать результат в память;
    \item изменить флаги.
\end{enumerate}

Такую команду трудно уложить в простую равномерную пятистадийную схему.

\subsection{Современный подход CISC}

Современные x86-процессоры решают проблему так:

\[
\text{сложные x86-команды}
\rightarrow
\text{декодирование}
\rightarrow
\text{микрооперации}
\rightarrow
\text{RISC-подобное ядро}.
\]

Внутренний конвейер работает преимущественно с микрооперациями фиксированного или более регулярного формата.

\subsection{Микро- и макрослияние}

В современных процессорах встречаются оптимизации:

\begin{itemize}
    \item \textbf{micro-op fusion} --- объединение нескольких микроопераций во внутренне более компактное представление;
    \item \textbf{macro-fusion} --- объединение некоторых пар машинных инструкций на стадии декодирования.
\end{itemize}

Эти методы уменьшают нагрузку на внутренние очереди и повышают пропускную способность.

\section{Глубокие конвейеры}

\begin{definition}
\textbf{Глубина конвейера} --- число стадий, на которые разбито выполнение инструкции.
\end{definition}

Увеличение глубины конвейера может позволить повысить тактовую частоту, потому что каждая стадия выполняет меньший объём работы.

Но глубокий конвейер имеет недостатки:

\begin{itemize}
    \item увеличивается штраф за ошибку предсказания перехода;
    \item возрастает сложность пересылки данных;
    \item увеличиваются накладные расходы на регистры конвейера;
    \item растёт энергопотребление;
    \item сложнее балансировать стадии.
\end{itemize}

Если длительность самой медленной стадии равна $t_{\max}$, а накладные расходы на конвейерный регистр равны $t_r$, то период такта примерно:

\[
    T_{\text{cycle}} \geq t_{\max} + t_r.
\]

Чрезмерное дробление стадий перестаёт давать выигрыш, потому что $t_r$ становится существенным.

\section{Связь конвейеризации, суперскалярности и внеочередного выполнения}

Эти методы решают разные, но связанные задачи.

\begin{itemize}
    \item \textbf{Конвейеризация} повышает пропускную способность за счёт перекрытия стадий разных инструкций.
    \item \textbf{Суперскалярность} позволяет обрабатывать несколько инструкций за такт.
    \item \textbf{Предсказание переходов} обеспечивает непрерывную выборку команд.
    \item \textbf{Спекулятивное выполнение} позволяет выполнять команды до окончательного подтверждения правильности пути.
    \item \textbf{Внеочередное выполнение} позволяет обходить задержки и использовать независимые инструкции.
    \item \textbf{Кэширование} снижает среднее время доступа к командам и данным.
\end{itemize}

Современный высокопроизводительный процессор обычно сочетает все эти методы.

\section{Обобщённая схема современного ядра}

Упрощённо современное ядро можно представить так:

\[
\begin{array}{c}
\text{Предсказатель переходов} \\
\downarrow \\
\text{Выборка команд} \\
\downarrow \\
\text{Декодирование} \\
\downarrow \\
\text{Переименование регистров} \\
\downarrow \\
\text{Очередь микроопераций / планировщик} \\
\downarrow \\
\text{Внеочередное выполнение} \\
\downarrow \\
\text{Буфер переупорядочивания} \\
\downarrow \\
\text{Commit в программном порядке}
\end{array}
\]

\section{Минимальный сквозной пример}

Рассмотрим фрагмент:

\begin{lstlisting}
1: LOAD R1, [A]
2: ADD  R2, R1, R3
3: SUB  R4, R5, R6
4: BEQ  R2, R0, L
5: MUL  R7, R8, R9
L:
6: XOR  R10, R11, R12
\end{lstlisting}

\subsection{Что происходит в простом конвейере}

\begin{itemize}
    \item Команда 2 зависит от команды 1 по RAW.
    \item Если данные из памяти не готовы, команда 2 ждёт.
    \item Команда 4 создаёт управляющий конфликт.
    \item Если переход предсказан неверно, команда 5 может быть ошибочно выбрана и затем отменена.
\end{itemize}

\subsection{Что происходит в современном процессоре}

\begin{enumerate}
    \item Предсказатель переходов пытается предсказать исход \texttt{BEQ}.
    \item Команды выбираются по предсказанному пути.
    \item Инструкции декодируются в микрооперации.
    \item Регистры переименовываются.
    \item \texttt{SUB} может выполниться независимо от \texttt{LOAD}.
    \item \texttt{ADD} ждёт результат загрузки.
    \item Если предсказание перехода верно, спекулятивная работа сохраняется.
    \item Если неверно, спекулятивные результаты отменяются через ROB.
\end{enumerate}

\section{Краткие выводы}

\begin{enumerate}
    \item CISC ориентирован на сложные команды и высокую плотность кода, но сложнее для декодирования и конвейеризации.
    \item RISC ориентирован на простые регулярные команды, что упрощает конвейеризацию и аппаратную реализацию.
    \item Конвейеризация повышает пропускную способность, но требует решения структурных, информационных и управляющих конфликтов.
    \item Предсказание переходов уменьшает простои из-за условных переходов.
    \item Спекулятивное выполнение повышает производительность, но требует механизмов отмены и точных исключений.
    \item Внеочередное выполнение позволяет использовать независимые инструкции, когда часть команд ждёт данные.
    \item Суперскалярность использует параллелизм уровня инструкций, но ограничена зависимостями и ресурсами.
    \item Кэширование команд и данных критически важно из-за большого разрыва между скоростью процессора и основной памяти.
    \item Современные CISC-процессоры часто внутри используют RISC-подобные микрооперации.
\end{enumerate}

\end{document}
